\chapter{Evaluation Method}
\label{ch:evaluation}

Training produced the behaviours described in Chapter~\ref{ch:development}, but it also showed
why a learning curve cannot be the final test. The training distribution changes with the
curriculum, actions are sampled stochastically, and a recent-success counter remembers only a
moving window. A convincing live trajectory can therefore coexist with a weak exported policy.
The evaluator was built to separate those two kinds of evidence.

\section{Freezing the test before observing the result}

Evaluation runs inside the same simulator with training disabled. It loads an exported ONNX
model, fixes the random seed and difficulty, records every episode and writes a machine-readable
summary. Success conditions come from the immutable session rather than from values adjusted
after evaluation begins.

The three tasks require different evidence. Fixed hover currently has no binary definition, so
it reports duration, fuel, planar error, three-dimensional and component velocities, tilt and
angular rate. Moving-target hover succeeds only after capturing the initial target and two
relocated targets, each within a \SI{35}{s} limit. Leg landing requires a stable four-foot
touchdown and additionally records first-contact speed, tilt, angular rate, impulse, rebound,
stable-hold time, fuel, feet outside the pad and structural strikes. Binary rates are accompanied
by the Wilson interval from Equation~\eqref{eq:wilson}.

For tracking and landing, the fixed protocol uses evaluation seed 20257 and five difficulty
bands,
\[
  d\in\{0,0.25,0.5,0.75,1.0\}.
\]
Each band contains 50 episodes, producing 250 trials per model. This arrangement was chosen to
show where performance changes. A single aggregate result could otherwise hide a policy that is
perfect on easy starts and fails completely at the curriculum boundary. Wind and injected
faults are disabled so that the evaluation measures the trained base task.

\section{A complete continuous-metric hover evaluation}

The earlier \texttt{HoverSimple2} diagnostic was stopped after four episodes and exposed a
large vertical drift. Its replacement, \texttt{HoverFixed}, started from fresh weights and was
continued through two immutable revisions to 20,078,171 PPO decision steps. The revised
single-engine interface has 29 observations and three continuous actions, with no separate
engine-enable action; minimum throttle is 25\%. Environment and trainer seeds are both 1.

The final ONNX export was evaluated deterministically with seed 20257 for all 250 planned
episodes. Each episode used the same \SI{60}{s} horizon, without curriculum bands, wind, grid
fins or RCS. All 250 episodes completed at the horizon, so the result provides a full distribution
of state and control metrics. The summary nevertheless declares
\texttt{successMetricDefined=false}. Horizon completion is consequently reported as a
termination outcome rather than interpreted as binary hover success.

\section{A complete tracking evaluation}

\texttt{HoverTrackSimple3} completed 30,000,039 steps. Its saved runtime state reported nonlinear
curriculum progress 0.9757 and linear progress 0.9070. Those values describe where training
finished; they are not evaluation outcomes. The final ONNX model was tested across all five
fixed bands, including exact $d=1$, even though training had not remained at that boundary.

Revision 3 records one sandbox engine, \SI{40000}{kg} starting fuel, no grid fins, RCS or wind,
and environment and trainer seeds equal to 1. The curriculum value during evaluation is supplied
by the evaluator rather than restored from runtime progress. This made the sharp boundary found
in Chapter~\ref{ch:results} observable instead of averaging it into easier trials.

\section{From the L10 baseline to the final landing test}

Landing was evaluated twice because L10 exposed a specific actuator problem that L11 was built
to address. L10 started from fresh weights and was continued through four immutable revisions to
79,629,752 steps. It reached 60.25\% effective curriculum difficulty and provided the first
complete 250-episode landing baseline. That evaluation finished before L11 began, preserving a
chronological comparison rather than selecting the better run after inspecting both.

L11 again started from fresh weights. Revision 1 reached 90M steps; revision 2 continued the
same policy to its final 100,000,086-step export. By then, the saved runtime state recorded both
linear and effective curriculum progress at 100\%, with a 92.9\% recent training-success rate.
The 90M checkpoint was evaluated before the continuation and is used only to show what the final
ten million steps changed. The official landing result uses the final \texttt{Rocket.onnx}.

Both L10 and L11 use the five-band protocol, seed 20257 and 50 episodes per band. Inference is
deterministic. Every success must satisfy the band-specific centre, velocity, tilt, angular-rate,
four-foot-support and hold-time limits stored in the session. The evaluator also preserves the
failure reason and post-contact actuator state.

The meaning of $d=1$ differs between the two runs. For L10 it is an out-of-distribution boundary
test because the curriculum reached only about 60\%. For L11 it is an in-range test because the
policy trained at full difficulty. The shared protocol makes the numerical comparison simple,
but this difference is necessary when interpreting why each model failed.

\section{Using video without replacing evidence}

Training and landing videos are useful because they show motion that a table cannot: engine
sequencing, final flare, foot contact and recovery. They should nevertheless be treated as
illustrations. A selected clip can overrepresent the most attractive trajectory. Any video used
in the presentation should therefore display the run ID, checkpoint, evaluation seed and
difficulty, while the complete CSV and JSON summaries remain the quantitative record. The
public repository should contain the plotting instructions and artifact layout needed to trace
those figures back to the frozen runs.
